\ChapterImageStar[cap:Marco-Teorico]{Marco Teórico}{./images/fondo.png}\label{cap:marcoTeorico}
\mbox{}\\
\noindent
Para fundamentar el diseño de la~\NDITI{}, este capítulo revisa cinco referentes que se complementan entre sí:~\UML, SysML, ArchiMate, el modelo por capas y~\TOGAF. Cada uno ofrece una perspectiva distinta sobre la representación de arquitecturas tecnológicas, y su análisis conjunto permite identificar qué aspectos pueden adoptarse, adaptarse o superarse en la propuesta.
\section{Lenguaje Unificado de Modelado (\UML)}
\noindent
\UML~es un lenguaje de propósito general para especificar, visualizar, construir y documentar sistemas de software y de negocio~\cite{UML01}, con capacidad además para la validación de sistemas complejos~\citep{Booch01}. En el contexto de la Infraestructura de TI, su uso resulta útil para representar dependencias lógicas, sin embargo autores como~\cite{Lovelle01}~afirman que existe una ausencia de soporte para elementos de arquitectura, teniendo que recurrir a extensiones para lograr representarse de manera adecuada.

\section{SysML}
\noindent
Siguiendo la línea conceptual anterior, SysML se define como un perfil de~\UML~2, el cual amplía las capacidades de UML para poder modelar arquitecturas de sistemas para aplicaciones de ingeniería~\cite{SYSML01}. Su uso ha sido ampliamente documentado, por ejemplo,~\cite{Thompson01}~muestran que SysML permite modelar sistemas de infraestructuras interdependientes mediante los bloques y diagramas disponibles en el lenguaje. Sin embargo en algunos contextos se ha encontrado a estos modelos como demasiado formales, haciendo que herramientas como la verificación de modelos sea rara vez utilizada debido al esfuerzo significativo que esta supone.

\section{ArchiMate}
Otro lenguaje de modelado relevante para este estudio es ArchiMate, el cual se define como un lenguaje de modelado abierto e independiente, estandarizado por The Open Group, cuya finalidad es proporcionar una representación uniforme para la Arquitectura Empresarial (AE) ofreciendo un enfoque arquitectónico integrado que no solo describe y visualiza los diferentes dominios de la arquitectura, sino también sus relaciones y dependencias subyacentes~\citep{Vicente01}. Una de las metas centrales de este lenguaje es lograr ofrecer un enfoque integrado que describa los diferentes dominios de arquitectura y sus relaciones, lo anterior logrado a través de una estructura compuesta de tres capas principales (Negocio, Aplicación y Tecnología)\citep{Patzwahl01}. 
Sin embargo Archimate presenta ciertas limitaciones en su uso que ya han sido documentadas, más específicamente~\cite{Nilus01}~en su publicación argumenta que el lenguaje presenta ambigüedad semántica que puede originar modelos inconsistentes o engañosos; una falta de soporte general para el modelado de datos e implementaciones variadas entre las herramientas que utilizan este estándar y que en consecuencia afectan considerablemente su interoperabilidad.

\section{Modelo por Capas}
\noindent
Adicionalmente la arquitectura por capas constituye un modelo de diseño que organiza un sistema en niveles funcionales independientes, donde cada nivel asume responsabilidades específicas~\citep{Spray2023}. Este enfoque resulta especialmente útil en el contexto del diseño de infraestructuras de \TI, ya que permite organizar los componentes y servicios tecnológicos en niveles claramente definidos, facilitando su análisis, documentación y posterior especificación. Este enfoque también facilita que cada capa pueda analizarse o modificarse de forma independiente, lo que simplifica tanto la documentación como la eventual actualización de la arquitectura.

\section{TOGAF (\textit{The Open Group Architecture Framework})}
\noindent
De igual manera, el marco~\TOGAF~se ha consolidado como uno de los referentes más empleados en la gestión y desarrollo de arquitecturas empresariales. Este enfoque metodológico ofrece una estructura sistemática que posibilita la alineación entre la estrategia organizacional y los procesos de~\TI{}, optimizando la toma de decisiones y el aprovechamiento de los recursos \citep{Mumtaza2025}. Su arquitectura por fases, que comprende la planificación, el diseño, la implementación y la validación, ofrece una estructura metodológica aplicable al proceso de especificación de infraestructuras de TI. Este enfoque además brinda flexibilidad y capacidad de adaptación para incorporar la notación~\NDITI~dentro de distintos escenarios académicos y organizacionales del programa de Ingeniería de Sistemas y Computación.


\section{Conclusión del Marco de Referencia}
\noindent
En síntesis, los lenguajes y marcos revisados —\UML{}, SysML, ArchiMate, la arquitectura por capas y~\TOGAF— aportan perspectivas complementarias para el diseño de la~\NDITI. En conjunto, los cinco referentes revisados justifican la elección de ArchiMate como base estructural de la~\NDITI{}, mientras que~\UML~y~\TOGAF~aportan mecanismos de extensión y una orientación metodológica complementaria. Esta combinación responde directamente a las necesidades identificadas en el programa, donde conviven distintos niveles de abstracción que una sola notación, tomada sin adaptación, no logra cubrir satisfactoriamente.